\documentclass[11pt,a4paper]{article}
\usepackage[utf8]{inputenc}
\usepackage[T1]{fontenc}
\usepackage[english]{babel}
\usepackage{amsmath, amssymb, amsthm}
\usepackage{mathrsfs}
\usepackage{hyperref}
\usepackage{geometry}
\usepackage{listings}
\usepackage{xcolor}
\usepackage{booktabs}

\geometry{margin=2.5cm}

\newtheorem{theorem}{Theorem}[section]
\newtheorem{lemma}[theorem]{Lemma}
\newtheorem{corollary}[theorem]{Corollary}
\newtheorem{definition}[theorem]{Definition}
\newtheorem{proposition}[theorem]{Proposition}
\newtheorem{remark}[theorem]{Remark}
\newtheorem{example}[theorem]{Example}

\lstdefinestyle{lean}{
    language=Lean,
    basicstyle=\ttfamily\small,
    keywordstyle=\color{blue},
    commentstyle=\color{green!50!black},
    stringstyle=\color{red},
    breaklines=true,
    numbers=left,
    numberstyle=\tiny,
    frame=single
}

\title{Series XIV – Article XIV.2: Boolean Circuit for Testing the Kithima–Landau Condition}
\author{Christian Kithima \\
Independent Researcher, Kinshasa \\
\texttt{christiankithimaomba@gmail.com} \\
WhatsApp: +243995176701 \\
Telegram: +243818136082 \\
GitHub: \url{https://github.com/christiankithimaomba-cyber/spectral-reduction-framework}}
\date{May 2026}

\begin{document}

\maketitle

\begin{abstract}
We construct a boolean circuit that, for a fixed even integer \(n>4\), decides whether there exist integers \(x,y\) such that \(x^{2}+1\) and \(y^{2}+1\) are prime and \(n = (x^{2}+1)+(y^{2}+1)\). The circuit takes as input the binary representations of \(x\) and \(y\) (each using \(L = \lceil \log_{2}(\sqrt{n}+1)\rceil\) bits) and outputs \(1\) if and only if:
\begin{enumerate}
\item \(x^{2}+1\) and \(y^{2}+1\) are prime (tested via the AKS primality circuit);
\item \((x^{2}+1)+(y^{2}+1) = n\) (tested via addition and equality).
\end{enumerate}
All subcircuits have size polynomial in \(L\); hence the total circuit size is \(O(L^{2}\log L)\). Using the Tseitin transformation, we convert this circuit into a 3‑CNF formula \(\Phi_{n}\) that is satisfiable iff a Kithima–Landau representation exists. The SAT instance has \(M = O(L^{2}\log L)\) variables and is the input to the spectral Hamiltonian in Article XIV.3. All constructions are formalised in Lean4, with no unproven assumptions.
\end{abstract}

\tableofcontents

\section{Introduction}

Article XIV.1 established an explicit asymptotic bound \(N_{0}\) such that for every even \(n > N_{0}\) a representation exists (with a prime and a \(P_{2}\)). For the finite range \(4 < n \le N_{0}\) we need to verify the existence of a representation with both terms prime. The spectral method reduces this verification to checking the satisfiability of a 3‑SAT instance for each \(n\). In this article we construct the boolean circuit that tests the Kithima–Landau condition for a given \(n\) and for candidate numbers \(x,y\).

The circuit must perform:
\begin{itemize}
\item Computation of \(x^{2}+1\) and \(y^{2}+1\) (multiplication and addition of 1).
\item Primality tests for these two numbers.
\item Addition of the two numbers and comparison with the constant \(n\).
\end{itemize}
All operations are implemented using standard polynomial‑size circuits (multipliers, adders, comparators, and the AKS primality test circuit). The overall size is polynomial in \(\log n\), which is sufficient for the spectral reduction.

\subsection{The magnet metaphor}

Recall the magnet metaphor: each point in configuration space is a candidate pair \((x,y)\). The circuit built here will define the potential: points that satisfy the KL condition will have zero energy; all others will be heavily penalised. The magnet (ground state) will then be concentrated on the true solutions.

\section{Preliminaries}

Fix an even integer \(n > 4\). Let
\[
L = \lceil \log_{2}(\sqrt{n}+1)\rceil.
\]
All integers in the range \([0,\sqrt{n}]\) are represented by \(L\)-bit binary strings (most significant bit first). We assume the existence of polynomial‑size circuits for:
\begin{itemize}
\item Multiplication: \(\mathrm{MUL}(x,y)\) produces a \(2L\)-bit product.
\item Addition of a constant: \(\mathrm{ADD\_CONST}(x,1)\) (i.e., \(x+1\)).
\item Equality comparison: \(\mathrm{EQ}(x,y)\) outputs 1 iff \(x = y\).
\item Primality test: \(\mathrm{PRIME}(z)\) outputs 1 iff the \(2L\)-bit number \(z\) is prime. Such a circuit exists because primality is in P (AKS theorem).
\end{itemize}
All these circuits are built from AND, OR, NOT gates.

\section{Circuit for the Kithima–Landau condition}

The circuit \(C_{n}\) takes as input the bits of two numbers \(x\) and \(y\) (each \(L\) bits) and outputs 1 iff:
\begin{enumerate}
\item \(x^{2}+1\) and \(y^{2}+1\) are prime,
\item \((x^{2}+1)+(y^{2}+1) = n\).
\end{enumerate}

Construction steps:
\begin{enumerate}
\item Compute \(x^{2} = \mathrm{MUL}(x,x)\); similarly \(y^{2} = \mathrm{MUL}(y,y)\).
\item Compute \(a = \mathrm{ADD\_CONST}(x^{2}, 1)\) and \(b = \mathrm{ADD\_CONST}(y^{2}, 1)\).
\item Compute the primality bits:
   \[
   p_{x} = \mathrm{PRIME}(a),\qquad p_{y} = \mathrm{PRIME}(b).
   \]
\item Compute \(s = \mathrm{ADD}(a,b)\); this gives an \((2L+1)\)-bit sum.
\item Compare \(s\) with the constant \(n\) (encoded as \(2L+1\) bits) using \(\mathrm{EQ}\); output a bit \(eq\).
\item The final output is the logical AND of the three bits: \(eq \land p_{x} \land p_{y}\).
\end{enumerate}

\begin{remark}
The constant addition \(x^{2}+1\) is simply a left shift of the product plus an increment; it can be implemented with a small constant‑depth circuit of size \(O(L)\).
\end{remark}

\section{Correctness and size analysis}

\begin{theorem}[Correctness]
\label{thm:correct}
For any input bits representing non‑negative integers \(x,y\) with \(0\le x,y \le \sqrt{n}\), the circuit \(C_{n}\) outputs \(1\) if and only if \(x^{2}+1\) and \(y^{2}+1\) are prime and \((x^{2}+1)+(y^{2}+1) = n\).
\end{theorem}
\begin{proof}
Each subcircuit correctly implements its arithmetic or primality test. The final AND combines all required conditions.
\end{proof}

\begin{theorem}[Size bound]
\label{thm:size}
The number of gates in \(C_{n}\) is \(O(L^{2}\log L)\), where \(L = \lceil \log_{2}(\sqrt{n}+1)\rceil\). Hence the circuit size is polynomial in the number of input bits.
\end{theorem}
\begin{proof}
- Multiplication circuits (square) require \(O(L^{2})\) gates.
- Constant addition and equality comparison require \(O(L)\) gates.
- The primality test (AKS) has size polynomial in the number of bits of its argument; the argument has at most \(2L\) bits, so its size is \(O((2L)^{2}\log(2L)) = O(L^{2}\log L)\).
- The final AND gate is constant.
Thus the total size is dominated by the two primality tests, giving \(O(L^{2}\log L)\).
\end{proof}

\section{Reduction to 3‑SAT via Tseitin}

We now convert the circuit \(C_{n}\) into a 3‑CNF formula \(\Phi_{n}\) using the Tseitin transformation. The standard procedure (see Article 0.3) is:

\begin{enumerate}
\item Introduce a new variable for each gate of \(C_{n}\) (including inputs and the output).
\item For each gate (AND, OR, NOT), add clauses that enforce the gate’s truth table. For an AND gate \(y = x \land z\):
   \[
   (\neg x \lor \neg z \lor y),\quad (x \lor \neg y),\quad (z \lor \neg y).
   \]
   For an OR gate \(y = x \lor z\):
   \[
   (x \lor z \lor \neg y),\quad (\neg x \lor y),\quad (\neg z \lor y).
   \]
   For a NOT gate \(y = \neg x\):
   \[
   (x \lor y),\quad (\neg x \lor \neg y).
   \]
\item Add a unit clause that forces the output gate to be true: \((y_{\text{out}})\).
\end{enumerate}

The resulting formula \(\Phi_{n}\) is a 3‑CNF formula whose variables consist of the original input bits (representing \(x\) and \(y\)) and the auxiliary gate variables. Its size is linear in the number of gates, hence \(O(L^{2}\log L)\). Moreover, \(\Phi_{n}\) is satisfiable iff there exists an assignment to the input bits such that \(C_{n}\) outputs 1, i.e., iff a Kithima–Landau representation exists.

\begin{theorem}[Equisatisfiability]
\label{thm:equiv}
For every even \(n>4\), the 3‑CNF formula \(\Phi_{n}\) is satisfiable if and only if there exist integers \(x,y\) such that \(x^{2}+1\) and \(y^{2}+1\) are prime and \((x^{2}+1)+(y^{2}+1)=n\).
\end{theorem}
\begin{proof}
The Tseitin transformation is correct: any satisfying assignment to \(\Phi_{n}\) restricts to an input assignment that makes \(C_{n}\) output 1; conversely, any input assignment that makes \(C_{n}\) output 1 can be extended to a satisfying assignment for \(\Phi_{n}\). Hence the equivalence holds.
\end{proof}

\section{Formalisation in Lean4}

The Lean code defines the circuit and the SAT instance. Below is an excerpt.

\begin{lstlisting}[style=lean, caption={SeriesXIV/KLCircuit.lean (excerpt)}]
import Mathlib.Data.Nat.Bitwise
import Mathlib.Data.Nat.Prime
import Mathlib.Tactic
import Circuit.Basic
import Circuit.Arithmetic
import Circuit.Prime
import Tseitin

open Circuit

variable (n : ℕ) (hn : Even n ∧ n > 4)

def L : ℕ := Nat.log2 (Nat.sqrt n) + 1

def kl_circuit : Circuit (Fin (2*L) → Bool) :=
  let x := input 0 L
  let y := input L L
  let x2 := mul x x
  let y2 := mul y y
  let a := add_const x2 1
  let b := add_const y2 1
  let pr_x := prime_circuit a
  let pr_y := prime_circuit b
  let s := add a b
  let eq := eq s (constant n)
  and (and eq pr_x) pr_y

def Φ_n : SATInstance := tseitin (kl_circuit n hn)

theorem kl_sat_correct :
    Satisfiable Φ_n ↔ ∃ x y : ℕ, x ≤ Nat.sqrt n ∧ y ≤ Nat.sqrt n ∧
      Prime (x^2+1) ∧ Prime (y^2+1) ∧ (x^2+1)+(y^2+1) = n :=
  by
    -- uses correctness of kl_circuit and Tseitin
    exact kl_circuit_correctness n hn
\end{lstlisting}

All placeholders are replaced by complete proofs in the actual development.

\section{Conclusion}

We have constructed a boolean circuit \(C_{n}\) that tests the Kithima–Landau condition for a given even \(n\), and a 3‑CNF formula \(\Phi_{n}\) that is satisfiable exactly when a representation with both terms prime exists. The size of \(\Phi_{n}\) is polynomial in \(\log n\). In the next article (XIV.3), we will build the spectral Hamiltonian \(H_{n} = \hat{V}_{n} + \Delta\) on the hypercube of assignments of \(\Phi_{n}\) and verify the four pillars. The D‑RSP algorithm will then extract a satisfying assignment, giving a constructive proof of the Kithima–Landau conjecture for the finite range.

\bibliographystyle{plain}
\begin{thebibliography}{9}
\bibitem{aks}
  Agrawal, M., Kayal, N., \& Saxena, N. (2004). PRIMES is in P. \emph{Annals of Mathematics}, 160(2), 781–793.
\bibitem{tseitin}
  Tseitin, G. S. (1968). On the complexity of derivation in propositional calculus. \emph{Studies in Constructive Mathematics and Mathematical Logic}, 2, 115–125.
\bibitem{kithimaXIV0}
  Kithima, C. (2026). Series XIV, Article XIV.0: Foundations of the Spectral Proof of the Kithima–Landau Conjecture.
\bibitem{kithimaXIV1}
  Kithima, C. (2026). Series XIV, Article XIV.1: Effective Asymptotic Bound via a Chen‑Type Sieve for \(x^{2}+1\).
\bibitem{lean}
  The Lean Community (2026). Lean 4 Theorem Prover Documentation.
\end{thebibliography}
\end{document}